\section{Related Work}
\label{relatedwork}

There are currently three dominant paradigms for large-scale pretraining of versatile image representations: self-supervised visual learning~\cite{chen2020simclr,assran2023ijepa,he2022mae}, image-text contrastive learning~\cite{radford2021clip,jia2021align,cherti2023reproducible} or image captioning~\cite{tschannen2023cappa,tschannen2024locca,fini2025aimv2}.
In this work, we build on top of self-supervised and image-text contrastive techniques to construct visual representations that are natively text-aligned and spatially aware.
We review relevant work in these areas and contextualize our advances with respect to them.

\noindent\textbf{Self-supervised learning} methods enable visual pretraining without the use of any labels, by formulating pretext tasks that derive supervisory signals from the images themselves.
Landmark approaches in this area, such as DINO~\cite{caron2021dino} and iBOT~\cite{zhou2022ibot}, enforce representational invariances to image augmentations such as crops, flips, color distortions and masking.
These methods adopt a student-teacher distillation setup, where the teacher model is constructed as an exponential moving average (EMA) of the student~\cite{he2020moco}.
DINOv2~\cite{oquab2024dinov2} combines these strategies into a well-tuned recipe which was scaled on a curated training set.
More recently, DINOv3~\cite{siméoni2025dinov3} and WebSSL~\cite{fan2025scalinglanguagefreevisualrepresentation} scale this recipe even further to image encoders with $7$B parameters, on $2$B curated images.

Our proposed \tipsvtwo framework leverages self-supervised losses, making two contributions to enhance them.
First, we extend the iBOT~\cite{zhou2022ibot} masked image modeling (MIM) strategy to introduce our novel iBOT++ technique, which enables direct supervision on both visible and masked tokens at pretraining time -- this simple change dramatically boosts patch-text alignment and benefits several downstream tasks.
Examples in the literature have some similarities to the iBOT++ strategy.
DMAE~\cite{wu2023dmae} predicts visible and masked patches (not tokens), focusing on robust classifiers.
SimMIM~\cite{xie2022simmim} also attempted predicting visible patches, but found it ineffective.
MaskAlign~\cite{xue2023maskalign} and MR-MAE~\cite{gao2023mrmae} align student visible tokens with those of pretrained frozen models.
To the best of our knowledge, overall, supervising both visible and masked tokens as part of MIM pretraining was unexplored.
Second, we introduce a significant simplification to the widely-used EMA learning setup~\cite{he2020moco}, where only projection heads undergo EMA updates and a single vision encoder is used.
This is particularly effective in our setup because the additional image-text contrastive loss generally provides a stable signal preventing model collapse.
This reduces the number of trainable parameters by nearly half, enabling much more efficient scaling.
Other methods which combine SSL and contrastive image-text techniques are SILC~\cite{naeem2024silc}, TIPS~\cite{tips_paper} and SigLIP2~\cite{tschannen2025siglip2multilingualvisionlanguage}: all of them employ EMA teachers but do not consider simplifications in this component.



\noindent\textbf{Contrastive vision-language pretraining} was introduced by CLIP~\cite{radford2021clip} and ALIGN~\cite{jia2021align}, leveraging web-scale image datasets with weakly supervised text labels.
More recent instantiations in this area include OpenCLIP~\citep{cherti2023reproducible}, SigLIP~\citep{zhai2023siglip}, Perception Encoder~\cite{bolya2025perceptionencoderbestvisual}, EVA~\citep{fang2023eva,fang2024eva02,sun2023evaclip} and MetaCLIP~\cite{xu2024metaclip,chuang2025metaclip2}.
Some methods in this area~\cite{lai2024veclip,fan2023laclip,tips_paper,stone2025visualcomposition} explore the use of synthetic image captions for contrastive image-text training.
In our \tipsvtwo model, we go beyond these to combine noisy web captions with synthetic captions at different granularities, providing a range of possible textual descriptions for images, increasing the robustness of the model.

Knowledge distillation~\cite{hinton2015distilling} techniques have been frequently used to train smaller vision-language encoders~\cite{tips_paper,tschannen2025siglip2multilingualvisionlanguage,bolya2025perceptionencoderbestvisual}, enabling high performance at a much cheaper compute budget.
In this work, one of our contributions is to show for the first time that distillation procedures, tuned to strengthen spatial awareness, can be helpful to enhance patch-text alignment of these models.
Surprisingly, even large pretrained models with weak patch-text alignment can be distilled into smaller models to enhance this capability.

















